% ===========================================================================
% abstract.tex — 中文摘要
% 对应格式要求：
%   - "摘要"字体为黑体小四号，水平居中
%   - 内容另起一行，字体为宋体小四号，行距18磅
%   - "关键词"另起一行，字体为黑体小四号
%   - 词条为宋体小四号，词条之间用分号隔开
%   - 关键词4~6个
%   - 摘要内容字数不少于800字（博士论文要求）
% ===========================================================================

\pagestyle{frontmatterstyle}

\begin{center}
  {\zihao{-4}\heiti 摘要}\\[6pt]
\end{center}

% 使用 \normalsize（小四号宋体，行距18磅）
\normalsize

% 中文摘要内容
% 【注意】博士论文中文摘要要求不少于800字，此处为示例模板，请按实际替换
水稻作为主要的粮食作物之一，其物候的精准监测与早期实现产量精准估测对保障粮食安全、优化农业生产管理和农业资源配置具有重要意义。传统的水稻物候监测依赖人工定点观测，存在耗时耗力、主观性强、覆盖范围有限等问题，难以实现大范围、高频次的动态监测。产量估测则多基于统计模型或卫星遥感数据，前者受限于样本代表性，后者存在空间分辨率低、易受云层遮挡等局限性。

无人机遥感技术凭借其高时空分辨率、低成本、机动灵活等优势，为田间尺度的作物监测提供了全新的技术手段。机器学习算法具备强大的数据挖掘能力、特征学习能力与非线性拟合能力，能够有效解析遥感数据与水稻物候、产量之间复杂的内在联系。因此，融合无人机遥感与机器学习技术，构建适配田间尺度的高效、精准的水稻物候监测识别与产量估测模型，对于提升水稻生产管理精细化水平、保障粮食安全具有重要的理论意义与应用价值。

本研究以水稻为研究对象，基于无人机多光谱遥感影像，系统开展了水稻物候期监测识别与产量估测研究。主要研究内容与结果如下：

（1）研究了基于无人机多光谱影像的水稻关键物候期判别方法。通过分析不同生育期水稻冠层光谱反射率的变化规律，提取了多个植被指数，结合随机森林、支持向量机等机器学习算法，构建了水稻物候期识别模型。结果表明，基于红边波段和近红外波段组合的植被指数在物候期判别中表现最佳，总体分类精度达到95\%以上。

（2）构建了基于多时相无人机遥感数据的水稻产量估测模型。利用水稻整个生长期的多光谱数据，提取了关键生育期的植被指数和纹理特征，分别采用偏最小二乘回归、随机森林回归、深度神经网络等方法建立了产量估测模型。研究表明，将抽穗期和灌浆期的特征进行融合，可以有效提高估测精度，R²达到0.85以上。

（3）提出了基于时序植被指数动态特征的水稻产量早期估测方法。利用水稻生长前期（移栽至抽穗期）的时序遥感数据，提取了植被指数曲线的动态特征参数，包括生长速率、峰值时间、曲线下面积等，建立了早期产量估测模型。该方法可在水稻收获前30\~{}45天实现产量的相对准确估测，为农业生产管理提供及时决策支持。

（4）开发了基于深度学习的端到端水稻产量估测方法。构建了基于卷积神经网络和长短期记忆网络的深度学习模型，直接从无人机多光谱影像序列中学习时空特征，避免了传统方法中人工特征提取的繁琐过程。实验结果表明，深度学习方法相较于传统机器学习方法，估测精度进一步提高，RMSE降低约15\%。

本研究构建的基于无人机遥感和机器学习的水稻物候监测与产量估测方法体系，可为精准农业中的作物生长监测和产量预测提供科学依据和技术支持，对保障国家粮食安全具有重要参考价值。

\vspace{0.5cm}

% 关键词
\noindent{\zihao{-4}\heiti 关键词：}%
{\zihao{-4}\songti 无人机遥感；机器学习；水稻物候监测；产量估测；深度学习}

\clearpage
